\documentclass[10pt,a4paper]{amsart}
\usepackage[margin=19mm]{geometry}
\usepackage{amsmath,amssymb,mathtools}
\usepackage[scr=rsfs,cal=euler]{mathalfa}
\usepackage{xfrac}
\usepackage[unicode,hidelinks,bookmarks=false]{hyperref}
\newcommand{\ie}{i.e.\ }
\newcommand{\cf}{cf.\ }
\newcommand{\resp}{respectively\ }
\newcommand{\Subsection}{\subsection}
\providecommand{\nobreakdash}{\nobreak\hbox{-}}
\setlength{\emergencystretch}{2em}
\multlinegap0pt
\usepackage{dsfont}
\usepackage{xcolor}
\usepackage{amssymb}
\usepackage{float}
\newtheorem{assumption}{Assumption}
\newtheorem{theorem}{Theorem}
 \newcommand{\E}{\mathbb{E}}            % expectation
 \newcommand{\Ll}{\langle}
 \newcommand{\PP}{\mathbb{P}}
 \newcommand{\R}{\mathbb{R}}\newcommand{\RR}{\mathbb{R}}
 \newcommand{\Rr}{\rangle}
\title{Self-normalized Cram\'er-type Moderate Deviation of Stochastic Gradient Langevin Dynamics}
\author{Hongsheng Dai, Xiequan Fan, Jianya Lu}
\date{Source: arXiv:2410.22047v1 (CC BY 4.0)}
\begin{document}
\maketitle

\noindent\emph{Source and licence.} Self-normalized Cram\'er-type Moderate Deviation of Stochastic Gradient Langevin Dynamics. Version arXiv:2410.22047v1 (CC BY 4.0), distributed by arXiv under the Creative Commons Attribution 4.0 International licence, http://creativecommons.org/licenses/by/4.0/, as recorded in the arXiv licence metadata for this exact version and shown on the abstract page https://arxiv.org/abs/2410.22047v1. Full licence notice: \texttt{licenses/arxiv-2410.22047-CC-BY-4.0.txt}.\par\smallskip

\noindent\textit{Editorial scope.} Counted unit: the article's theorem coro:berry (source0.tex lines 426--435). The article's complete original proof of it is delivered with it (source0.tex lines 966--1014, one \texttt{proof} environment of 49 lines, which the article places away from the statement). No mathematics is written, repaired or re-derived: the statement and the proof are the article's own lines, in the article's own notation, and no retained line is substituted or re-flowed.  Retained uncounted as the source-local context the proof needs: lines 350--358; lines 362--367; lines 401--421.  Nothing was omitted from any retained span, so the package carries no omission record.  No retained line contains a citation, so the wrapper carries no bibliography and no bibliography entry.  No retained line contains a figure environment, an image-inclusion command or drawing code, so no image is delivered and the package declares no asset dependency.  Editorial formatting only, in addition: an AMS \texttt{amsart} preamble that restates the article's own declarations and macros used by the retained lines, namely the environments \texttt{assumption}, \texttt{theorem} on separate counters, together with the standard packages those lines need.  The retained environments print in this document as Assumption 1, Theorem 1 in order of appearance, whereas the article prints the same items with the numbering of its own full document.  The complete native source of the article as distributed in the arXiv e-print for this exact version is bundled unchanged as \texttt{sources/167591/source0.tex}.
\begin{assumption}\label{assu1}
	There exist constants $L, K_1 >0$ and $K_2 \ge 0$ such that for every $x,y\in \R^d$, $z\in\R^r$, 
	\begin{eqnarray}\label{e:lip}
		|\nabla\psi(x,z)-\nabla\psi(y,z)|\le L|x-y|,
	\end{eqnarray}
	\begin{eqnarray}\label{e:dissi}
		\Ll x-y,-\nabla \psi(x,z)+\nabla \psi(y,z)\Rr\le -K_1|x-y|^2+K_2.
	\end{eqnarray}
\end{assumption}

\begin{assumption}\label{assu2}
	The random variable $\nabla\psi(x,\zeta)$ is sub-Gaussian for any $x\in\R^d$, that is, there exist positive constants $K_\zeta$ and $C$ such that 
	\begin{eqnarray*}
		\E\exp\{K_\zeta|\nabla\psi(x,\zeta)|^2\} \le C.
	\end{eqnarray*} 
\end{assumption}

\begin{theorem}\label{thm-self-normal}
Suppose Assumptions \ref{assu1} and \ref{assu2} hold. Let $\omega_0\sim\pi_\eta$ and $h \in \mathrm{Lip}_1(\R^d,\R)$, set $m=o(\eta^{-2})$ with $m\eta\to\infty$. Then we have:

(i) as $m\le\eta^{-\frac{13}8}\delta^{-\frac98}$,
\begin{eqnarray*}
	\Big| \frac{\PP(\mathcal W_\eta> x)}{1-\Phi(x)}
	-1 \Big| & \le & 
	C\big(x^3m^{-\frac14}+(1+x)m^{-\frac14}\ln m+x^6\eta^{\frac12}\delta^{-\frac12}\big)
\end{eqnarray*}
uniformly for $0\le x=o(\eta^{{-\frac{1}{12}}}\delta^{\frac{1}{12}})$ as $\eta$ tends to zero and $m$ tends to infinity.  

(ii) as $m>\eta^{-\frac{13}8}\delta^{-\frac98}$,
\begin{eqnarray*}
	\Big| \frac{\PP(\mathcal W_\eta> x)}{1-\Phi(x)}
	-1 \Big| & \le & 
	C\big(x^3(m\eta\delta)^{-\frac12}+\sqrt m\eta\delta x+ m^{-\frac14}\ln m\big)
\end{eqnarray*}
uniformly for $0\le x=o( (m\eta\delta)^{\frac16}\wedge (\sqrt m\eta\delta)^{-1})$ as $\eta$ tends to zero and $m$ tends to infinity. 

Moreover, the same results hold when $\mathcal{W}_\eta$ is replaced by $-\mathcal{W}_\eta$.
\end{theorem}	

\begin{theorem}\label{coro:berry}
Under the assumption of Theorem \ref{thm-self-normal}, we have
\begin{eqnarray*}
	\sup_{x\in\R}|\PP(\mathcal{W}_\eta\le x)-\Phi(x)|\le Cm^{-\frac14}\ln m.
\end{eqnarray*}	
Further assume $m=C\eta^{-2}/|\ln\eta|$, we have
\begin{eqnarray*}
	\sup_{x\in\R}|\PP(\mathcal{W}_\eta\le x)-\Phi(x)|\le C\eta^{\frac12}|\ln \eta|^{5/4}.
\end{eqnarray*}	
\end{theorem}

\begin{proof}[Proof of Theorem \ref{coro:berry}]
For the case $m\le\eta^{-\frac{13}8}\delta^{-\frac98}$, denote $C_{m,\eta}=\eta^{{-\frac{1}{24}}}\delta^{\frac{1}{24}} $. It is easy to have the following decomposition,	
\begin{eqnarray*}
	\sup_{x\in\R}|\PP(\mathcal{W}_\eta< x)-\Phi(x)|&\le& 	\sup_{x\le -C_{m,\eta}}|\PP(\mathcal{W}_\eta\le x)-\Phi(x)|+\sup_{-C_{m,\eta}\le x\le 0 }|\PP(\mathcal{W}_\eta\le x)-\Phi(x)|\\
	&&+\sup_{0\le x\le C_{m,\eta}}|\PP(\mathcal{W}_\eta\le x)-\Phi(x)|+\sup_{x> C_{m,\eta}}|\PP(\mathcal{W}_\eta\le x)-\Phi(x)|\\
	&=:& I_1+I_2+I_3+I_4.
\end{eqnarray*}		


For $I_1$ and $I_4$,  Theorem \ref{thm-self-normal} implies
\begin{eqnarray*}
I_1&=& \sup_{x\le -C_{m,\eta}}|\PP(\mathcal{W}_\eta\le x)-\Phi(x)|\\
&\le& \sup_{x\le -C_{m,\eta}}\PP(\mathcal{W}_\eta\le x)+\Phi(-c_{\eta,m})\\
&\le& \Phi(-C_{m,\eta})e^C+\Phi(-C_{m,\eta})\le Cm^{-\frac14}\ln m.
\end{eqnarray*}
Similarly,
\begin{eqnarray*}
I_4\le Cm^{-\frac14}\ln m.
\end{eqnarray*}
For $I_2$ and $I_3$, Theorem \ref{thm-self-normal} and the inequality $|e^x-1|\le|x|e^{|x|}$ imply
	\begin{eqnarray*}
	I_2&=&\sup_{-C_{m,\eta}\le x\le 0 }|\PP(\mathcal{W}_\eta\le x)-\Phi(x)|\\
	&\le&\sup_{-C_{m,\eta}\le x\le 0 }C\Phi(x)\big(x^3m^{-\frac14}+(1+x)m^{-\frac14}\ln m+x^6\eta^{\frac12}\delta^{-\frac12}\big)\\
	&\le& Cm^{-\frac14}\ln m.
	\end{eqnarray*}
	Similarly,
\begin{eqnarray*}
		I_3\le Cm^{-\frac14}\ln m.
\end{eqnarray*}
Combining the estimation for the terms $I_1$-$I_4$, we have
\begin{eqnarray*}
	\sup_{x\in\R}|\PP(\mathcal{W}_\eta< x)-\Phi(x)|&\le& 	Cm^{-\frac14}\ln m.
\end{eqnarray*}		

For the case $m>\eta^{-\frac{13}8}\delta^{-\frac98}$, taking $C_{m,\eta}= (m\eta\delta)^{\frac1{12}}\wedge (\sqrt m\eta\delta)^{-\frac12}$ instead of $m^{1/24}$, we can similarly show that
\begin{eqnarray*}
	\sup_{x\in\R}|\PP(\mathcal{W}_\eta< x)-\Phi(x)|&\le& Cm^{-\frac14}\ln m.
\end{eqnarray*}		

Thus, for any $\eta^{-1}<m=o(\eta^{-2})$, we have
\begin{eqnarray*}
	\sup_{x\in\R}|\PP(\mathcal{W}_\eta< x)-\Phi(x)|&\le& Cm^{-\frac14}\ln m.
\end{eqnarray*}		
Further assume $m=C\eta^{-2}/|\ln\eta|$, we can get
\begin{eqnarray*}
	\sup_{x\in\R}|\PP(\mathcal{W}_\eta\le x)-\Phi(x)|
	&\le& C\eta^{\frac12}|\ln \eta|^{5/4}.
\end{eqnarray*}	
\end{proof}	

\end{document}
